\chapter{Implementarea aplicației}
\label{chapter:implementation}

Implementarea urmează modulele definite în Capitolul \ref{chapter:application}. Limbajul principal este Python, folosit pentru scraping, procesarea datelor, antrenare și backend. Interfața este realizată cu HTML, CSS și JavaScript fără framework frontend, pentru a păstra deployment-ul compact. Codul este organizat în pachetul \texttt{carvaluator\_scraper}, iar comenzile de reproducere sunt prezentate în Anexa \ref{appendix:commands}.

\chapter{Colectarea și pregătirea datelor}
\label{chapter:data}

\section{Sursele de date}
\label{sec:data-sources}

Datele proiectului provin din anunțuri publice de pe \autovit\ și \mobilede. Alegerea acestor două surse are o justificare practică. \autovit\ este relevant pentru piața românească, iar \mobilede\ este una dintre platformele europene importante pentru autoturisme second-hand, inclusiv pentru cumpărătorii români care caută mașini din Germania.

Setul combinat folosit pentru evaluarea finală conține 9633 de rânduri păstrate după normalizare și deduplicare. Din acestea, 7057 provin din \autovit, iar 2576 din \mobilede. Creșterea volumului \mobilede\ a fost obținută printr-o metodă experimentală bazată pe pagini SEO convertite în Markdown, descrisă în secțiunea dedicată acestei surse.

\labelindexref{Figura}{fig:dataset-sources} prezintă distribuția surselor. Setul rămâne dominat de \autovit, dar ponderea \mobilede\ este suficient de mare pentru a contribui vizibil la distribuțiile de preț și la antrenarea finală.

\fig[width=0.70\textwidth]{src/img/carvaluator/dataset_sources.png}{fig:dataset-sources}{Componența setului de date după sursă}

\begin{table}[htb]
    \centering
    \caption{Rezumatul setului combinat înainte de curățarea finală pentru antrenare}
    \label{tab:data-summary}
    \begin{tabular}{lr}
        \toprule
        Indicator & Valoare \\
        \midrule
        Rânduri încărcate & 10137 \\
        Rânduri păstrate după deduplicare & 9633 \\
        Duplicate exacte eliminate & 291 \\
        Duplicate fuzzy eliminate & 213 \\
        Rânduri fără preț & 9 \\
        Mărci distincte & 82 \\
        Modele distincte & 693 \\
        Mediană preț brut & 16640 EUR \\
        Mediană kilometraj brut & 112927 km \\
        \bottomrule
    \end{tabular}
\end{table}

\section{Scraping pentru Autovit}
\label{sec:data-autovit}

Pentru \autovit, scraperul a fost construit pornind de la paginile de căutare. Fiecare pagină conține mai multe anunțuri, iar structura HTML include informații suficiente pentru extragerea linkului, titlului, prețului, anului, kilometrajului și câtorva câmpuri tehnice. În unele cazuri, datele sunt disponibile și în structuri JSON incluse în pagină, ceea ce permite o extragere mai stabilă decât parsarea textului vizibil.

Fluxul de scraping salvează fiecare anunț ca obiect JSONL. Alegerea JSONL este intenționată: fiecare linie este un document independent, iar fișierul poate fi procesat incremental. Dacă scraping-ul se oprește la pagina 80 din 150, datele deja salvate rămân valide. În plus, JSONL păstrează și câmpul \texttt{raw}, unde pot fi arhivate detalii specifice site-ului. Această decizie a fost utilă atunci când unele câmpuri au trebuit reinterpretate fără a reporni colectarea.

În practică, \autovit\ a ridicat mai multe dificultăți. Prima dificultate a fost instabilitatea structurii paginilor. Numele claselor HTML și poziția câmpurilor se pot modifica fără avertisment, iar scraperul trebuie să fie tolerant la câmpuri lipsă. A doua dificultate a fost deduplicarea. Anunțurile promovate sau republicate pot apărea pe mai multe pagini, iar identificatorul de listing nu este întotdeauna suficient pentru eliminarea tuturor cazurilor similare. A treia dificultate a fost legată de imagini. Unele imagini sunt încărcate lazy, iar alte imagini pot reprezenta elemente de dealer, nu mașina propriu-zisă.

\section{Scraping pentru mobile.de}
\label{sec:data-mobilede}

Integrarea \mobilede\ a fost mai dificilă decât integrarea \autovit. Site-ul are localizare pe mai multe limbi, pagini dinamice, protecții anti-bot și comportament diferit între paginile de căutare și paginile de detaliu. În multe încercări, cererile automate au returnat erori 401, 403 sau 429. Aceste coduri indică lipsă de autorizare, interdicție de acces sau limitare de rată.

Pentru a crește șansele de colectare, scraperul \mobilede\ încearcă să folosească endpoint-uri JSON cu headere de browser. În plus, există un fallback cu Playwright, util atunci când pagina trebuie randată într-un browser real. Totuși, fallback-ul nu garantează succesul, deoarece protecțiile pot detecta mediile automate sau pot afișa pagini de consimțământ și acces restricționat.

O problemă importantă a apărut și la imaginile anunțurilor \mobilede. Inițial, aplicația putea afișa imaginea greșită, de exemplu bannerul dealerului sau fotografia de showroom, nu fotografia principală a mașinii. Soluția a fost prioritizarea imaginilor din galeria anunțului și filtrarea surselor care par logo-uri sau imagini de dealer. Această problemă arată că scraping-ul nu înseamnă doar extragerea unui câmp, ci și interpretarea corectă a contextului în care apare acel câmp.

Din cauza acestor dificultăți, primele date \mobilede\ colectate au fost puține. Integrarea a rămas însă valoroasă pentru aplicație deoarece permite analizarea unor linkuri de detaliu și expune limitele reale ale unui sistem care depinde de site-uri externe.

Într-o etapă ulterioară, am testat și o strategie alternativă pentru mărirea volumului de date \mobilede. API-ul oficial de căutare \mobilede\ există, dar accesul la anunțuri necesită autentificare Basic și un cont cu drepturi pentru serviciul respectiv \cite{mobilede_search_api}. Fără aceste credențiale, endpoint-ul de căutare a răspuns cu 401, în timp ce endpoint-urile publice de referință au returnat doar liste de mărci și categorii, nu anunțuri concrete.

Metoda care a funcționat practic a fost citirea paginilor SEO publice de forma \url{https://www.mobile.de/en/car/marca/model} printr-un serviciu de conversie HTML în Markdown, Jina Reader \cite{jina_reader}. În loc să controleze un browser, scriptul descarcă reprezentarea Markdown a paginii, apoi parsează cardurile de anunțuri. Din fiecare card se pot extrage titlul, prețul, anul primei înmatriculări, kilometrajul, puterea, combustibilul, orașul și prima imagine. Pentru cilindree, cardurile nu oferă mereu un câmp separat, astfel că parserul folosește o regulă conservatoare din titlu, de exemplu \texttt{2.0 TDI}, \texttt{1.6 HDI} sau \texttt{1.5 TSI}. Regula evită codurile comerciale precum \texttt{35 TFSI} și expresiile de consum precum \texttt{5,5 l/100km}. Pentru că aceste carduri nu expun întotdeauna ID-ul real al anunțului, scriptul generează un identificator sintetic pe baza unei amprente formate din titlu, marcă, model, an, kilometraj, putere, preț și oraș.

Prin această metodă au fost citite 118 pagini de marcă sau model și au fost obținute 2735 rânduri brute \mobilede, după filtrarea blocurilor care nu reprezentau anunțuri. După exportul prin pipeline-ul de normalizare și deduplicare fuzzy au rămas 2534 rânduri curate din metoda Markdown. Combinarea cu lotul \mobilede\ obținut anterior a produs 2576 rânduri curate. În acest set \mobilede\ combinat, puterea este disponibilă pentru 2533 rânduri, iar capacitatea cilindrică pentru 647 rânduri. Această metodă nu înlocuiește API-ul oficial și nu este ideală pentru producție, deoarece depinde de un serviciu intermediar și nu oferă toate câmpurile tehnice, precum transmisia sau caroseria. Totuși, pentru un prototip academic, ea oferă un snapshot suplimentar util și reproductibil, salvat local în JSONL și CSV.

Pentru ca aceste rânduri să fie folosite la antrenare, filtrarea finală a fost ajustată. Valorile numerice imposibile sunt eliminate în continuare, însă câmpurile tehnice lipsă, precum capacitatea cilindrică, nu mai duc automat la eliminarea rândului. Ele sunt imputate ulterior de pipeline-ul de preprocesare. Această schimbare este importantă deoarece paginile Markdown oferă preț, an, kilometraj și putere, dar nu întotdeauna detalii complete despre motor.

\section{Schema datelor brute}
\label{sec:data-schema}

Fiecare rând brut conține câmpuri comune și câmpuri specifice sursei. Câmpurile comune sunt folosite ulterior pentru normalizare și modelare. \labelindexref{Tabelul}{tab:raw-fields} prezintă câteva exemple.

\begin{table}[htb]
    \centering
    \caption{Câmpuri extrase din anunțuri}
    \label{tab:raw-fields}
    \begin{tabular}{lp{8cm}}
        \toprule
        Câmp & Rol \\
        \midrule
        \texttt{source} & Platforma de proveniență, de exemplu \texttt{autovit} sau \texttt{mobilede} \\
        \texttt{listing\_id} & Identificatorul anunțului, atunci când poate fi extras \\
        \texttt{url} & Linkul original al anunțului \\
        \texttt{title} & Titlul afișat în anunț \\
        \texttt{make}, \texttt{model}, \texttt{version} & Marca, modelul și versiunea vehiculului \\
        \texttt{year}, \texttt{mileage\_km} & Anul și kilometrajul \\
        \texttt{price\_eur} & Prețul convertit sau standardizat în EUR \\
        \texttt{fuel\_type}, \texttt{transmission} & Câmpuri tehnice categorice \\
        \texttt{power\_hp}, \texttt{engine\_capacity\_cm3} & Putere și capacitate cilindrică \\
        \texttt{seller\_type}, \texttt{location\_city} & Context comercial și geografic \\
        \bottomrule
    \end{tabular}
\end{table}

\section{Normalizare}
\label{sec:data-normalization}

Datele brute sunt utile pentru audit, dar nu pot fi folosite direct în modele. Prețurile pot fi scrise cu separatori diferiți, kilometrajul poate conține text, combustibilul poate avea variante lingvistice, iar unele câmpuri pot lipsi. Normalizarea transformă aceste reprezentări într-un format comun.

Prețul este convertit în \texttt{price\_eur}. În datele curente, majoritatea anunțurilor sunt deja în EUR, dar câmpul păstrează explicit moneda originală pentru extensii viitoare. Kilometrajul este extras ca număr întreg în kilometri. Puterea este standardizată în cai putere, iar capacitatea cilindrică în centimetri cubi. Pentru combustibil, valorile sunt mapate în categorii precum \texttt{diesel}, \texttt{petrol}, \texttt{hybrid}, \texttt{plug\_in\_hybrid} și \texttt{electric}.

Normalizarea construiește și câmpuri auxiliare. \texttt{dedupe\_exact\_key} identifică duplicatele cu același ID de sursă. \texttt{dedupe\_fuzzy\_key} combină marca, modelul, anul, kilometrajul, prețul și vânzătorul într-o amprentă conservatoare. \texttt{completeness\_score} numără câte câmpuri relevante sunt prezente, fiind util pentru filtrarea rândurilor slabe.

\section{Deduplicare}
\label{sec:data-dedup}

Deduplicarea are două niveluri. Primul nivel elimină duplicate exacte, folosind cheia de sursă și identificatorul anunțului. Al doilea nivel elimină duplicate probabile. Acesta este necesar deoarece același vehicul poate apărea cu identificatori diferiți sau poate fi republicat.

Cheia fuzzy este construită astfel încât să fie conservatoare. Dacă două anunțuri au aceeași marcă, același model, același an, kilometraj foarte asemănător, preț foarte asemănător și același vânzător sau oraș, este probabil să fie același vehicul. În setul combinat extins, acest pas a eliminat 213 rânduri suplimentare după eliminarea duplicatelor exacte. Numărul a crescut față de experimentul inițial deoarece paginile SEO \mobilede\ pot afișa același anunț în mai multe pagini de marcă sau model.

\section{Analiza exploratorie}
\label{sec:data-eda}

Analiza exploratorie arată că datele sunt neuniforme. Distribuția prețurilor este asimetrică, cu multe anunțuri în zona 5000-25000 EUR și câteva vehicule premium cu valori mult mai mari. \labelindexref{Figura}{fig:price-distribution} ilustrează această distribuție.

\fig[width=0.78\textwidth]{src/img/carvaluator/price_distribution.png}{fig:price-distribution}{Distribuția prețurilor în setul de date}

Kilometrajul și prețul au o relație negativă, dar relația nu este perfectă. Mașinile noi cu kilometraj mic sunt mai scumpe, însă marca și segmentul pot schimba semnificativ poziția unui anunț. O mașină premium cu kilometraj mare poate rămâne mai scumpă decât o mașină de oraș cu kilometraj mic. \labelindexref{Figura}{fig:price-mileage} arată această relație.

\fig[width=0.82\textwidth]{src/img/carvaluator/price_vs_mileage.png}{fig:price-mileage}{Prețul în funcție de kilometraj}

\labelindexref{Figura}{fig:median-year} arată evoluția medianei prețului în funcție de anul mașinii. Tendința generală este crescătoare pentru mașinile mai noi, dar există variații cauzate de compoziția pe mărci și modele.

\fig[width=0.82\textwidth]{src/img/carvaluator/median_price_by_year.png}{fig:median-year}{Mediana prețului în funcție de anul mașinii}

Cele mai frecvente mărci sunt prezentate în \labelindexref{Figura}{fig:top-makes}. Distribuția confirmă faptul că setul de date reflectă preferințele pieței locale, cu mărci comune în România și Europa.

\fig[width=0.82\textwidth]{src/img/carvaluator/top_makes.png}{fig:top-makes}{Cele mai frecvente mărci în setul de date}

\section{Pregătirea pentru antrenare}
\label{sec:data-training-ready}

Setul final pentru antrenare nu include toate rândurile normalizate. Înainte de împărțirea train-test, pipeline-ul elimină exemple cu prețuri nerealiste, ani foarte vechi sau viitori nejustificați și valori tehnice extreme. Pentru antrenarea finală pe setul extins, raportul indică 9367 rânduri curate, dintre care 7493 au fost folosite la antrenare și 1874 la testare.

Caracteristicile candidate includ marca, modelul, combinația marcă-model, versiunea, anul, vârsta vehiculului, luna și anul primei înmatriculări, kilometrajul, kilometrajul pe an, combustibilul, transmisia, caroseria, puterea, capacitatea cilindrică, raportul cai putere pe litru, tipul vânzătorului și localizarea. Nu toate aceste caracteristici sunt păstrate. Selecția statistică este discutată în capitolul următor.


\chapter{Modelare și selecția caracteristicilor}
\label{chapter:models}

\section{Caracteristici derivate}
\label{sec:models-features}

Modelele nu folosesc doar câmpurile extrase direct din anunțuri. Pipeline-ul construiește și caracteristici derivate care surprind relații mai utile pentru predicție. Vârsta vehiculului este calculată ca diferență între anul curent și anul mașinii. Kilometrajul pe an împarte kilometrajul total la vârsta vehiculului, evitând interpretarea identică a două mașini cu același kilometraj, dar cu vârste diferite. Raportul cai putere pe litru este calculat din putere și capacitate cilindrică, oferind un indicator simplu al performanței specifice.

O caracteristică importantă este \texttt{make\_model}, combinația dintre marcă și model. Aceasta ajută modelele să învețe că prețul nu depinde independent de marcă și model, ci de perechea lor. De exemplu, marca BMW are mai multe modele cu prețuri foarte diferite, iar modelul Seria 3 are altă distribuție decât X5. O reprezentare combinată reduce ambiguitatea.

\begin{table}[htb]
    \centering
    \caption{Exemple de caracteristici derivate}
    \label{tab:derived-features}
    \begin{tabular}{lp{8cm}}
        \toprule
        Caracteristică & Interpretare \\
        \midrule
        \texttt{vehicle\_age} & Vechimea mașinii în ani \\
        \texttt{mileage\_per\_year} & Ritmul de utilizare anuală \\
        \texttt{hp\_per\_liter} & Putere raportată la capacitatea cilindrică \\
        \texttt{make\_model} & Segment comercial aproximat prin marcă și model \\
        \bottomrule
    \end{tabular}
\end{table}

\section{Selecția statistică}
\label{sec:models-selection}

Selecția caracteristicilor are rolul de a reduce zgomotul și de a elimina coloanele slab reprezentate. Pentru caracteristicile numerice, pipeline-ul calculează corelația Pearson față de prețul logaritmic. Pentru caracteristicile categorice, prețul este împărțit în intervale, iar asocierea dintre categorie și intervalul de preț este testată prin chi-pătrat. Mărimea efectului este exprimată printr-o valoare de tip Cramér V.

În configurația finală, pragurile au fost: \(\alpha = 0.05\), corelație Pearson minimă absolută \(0.05\) pentru numerice și Cramér V minim \(0.25\) pentru categorice. Aceste praguri sunt suficient de permisive pentru a păstra semnale slabe, dar elimină câmpuri cu puține valori sau cu asociere insuficientă.

\labelindexref{Figura}{fig:feature-significance} prezintă cele mai relevante caracteristici. Printre cele păstrate se află versiunea, puterea, anul, vârsta vehiculului, combinația marcă-model, kilometrajul, capacitatea cilindrică, orașul și anul primei înmatriculări.

\fig[width=0.88\textwidth]{src/img/carvaluator/feature_significance_top.png}{fig:feature-significance}{Caracteristici relevante după testele statistice}

\begin{table}[htb]
    \centering
    \caption{Caracteristici selectate în antrenarea finală}
    \label{tab:selected-features}
    \begin{tabular}{@{}p{3.2cm}p{9.5cm}@{}}
        \toprule
        Tip & Caracteristici \\
        \midrule
        Numerice & \texttt{year}, \texttt{vehicle\_age}, \texttt{mileage\_km}, \texttt{power\_hp}, \texttt{engine\_capacity\_cm3} \\
        Numerice derivate & \texttt{mileage\_per\_year}, \texttt{hp\_per\_liter} \\
        Categorice & \texttt{make}, \texttt{model}, \texttt{make\_model}, \texttt{version}, \texttt{seller\_type}, \texttt{location\_city} \\
        Înmatriculare & \texttt{first\_registration\_year} \\
        \bottomrule
    \end{tabular}
\end{table}

Unele câmpuri au fost eliminate deși pot fi importante în realitate. \texttt{transmission} și \texttt{body\_type} au fost eliminate în setul combinat deoarece erau completate pentru prea puține rânduri. \texttt{fuel\_type} a avut semnificație statistică, dar efectul a fost sub pragul ales. Această situație indică o limitare a datelor, nu neapărat o lipsă de relevanță a combustibilului.

\section{Preprocesare}
\label{sec:models-preprocessing}

Pipeline-ul folosește transformări diferite pentru variabile numerice și categorice. Variabilele numerice sunt imputate atunci când lipsesc și scalate pentru modelele sensibile la distanță, precum SVR și KNN. Variabilele categorice sunt imputate cu o valoare specială pentru lipsă și transformate prin one-hot encoding.

Această separare este importantă deoarece modelele au cerințe diferite. Arborii de decizie sunt mai puțin sensibili la scalare, dar SVR și KNN pot fi dominate de variabile cu valori mari, de exemplu kilometrajul. Prin standardizare, fiecare variabilă numerică are o contribuție comparabilă în spațiul de intrare.

\section{Modele candidate}
\label{sec:models-candidates}

Au fost antrenate șapte modele:

\begin{itemize}
    \item \texttt{svr\_rbf}, un SVR cu kernel RBF;
    \item \texttt{ridge}, o regresie liniară regularizată;
    \item \texttt{knn\_distance}, KNN cu ponderare după distanță;
    \item \texttt{random\_forest}, ansamblu Random Forest;
    \item \texttt{extra\_trees}, ansamblu Extra Trees;
    \item \texttt{gradient\_boosting}, boosting pe arbori;
    \item \texttt{voting\_ensemble}, ansamblu care combină mai multe modele.
\end{itemize}

Implementarea este realizată cu scikit-learn \cite{pedregosa2011scikit}. Fiecare model este inclus într-un pipeline care conține preprocesarea corespunzătoare. Această alegere reduce riscul de inconsistență între antrenare și predicție, deoarece aceeași transformare este aplicată și datelor din anunțurile live.

\section{Funcționarea modelelor folosite}
\label{sec:models-how-they-work}

Modelele au fost alese astfel încât să acopere familii diferite de regresie: modele liniare, modele bazate pe vecini, modele cu vectori suport și ansambluri de arbori. \labelindexref{Tabelul}{tab:model-methods} sintetizează ideea fiecărui algoritm și modul în care se interpretează în contextul evaluării unui anunț auto.

\begin{table}[htb]
    \centering
    \caption{Modul de funcționare al modelelor utilizate}
    \label{tab:model-methods}
    \begin{tabular}{@{}p{2.4cm}p{5.8cm}p{3.8cm}@{}}
        \toprule
        Model & Idee de funcționare & Exemplu în proiect \\
        \midrule
        SVR RBF & Caută o funcție neliniară care păstrează erorile mici într-un interval de toleranță și penalizează abaterile mari. Kernelul RBF permite relații neliniare între an, kilometraj, putere și preț. & Două mașini cu același an pot primi estimări diferite dacă distanța combinată față de exemplele de antrenare diferă mult. \\
        Ridge & Este o regresie liniară cu regularizare. Coeficienții sunt micșorați pentru a evita supraînvățarea pe multe coloane one-hot. & O marcă sau o versiune poate avea un coeficient pozitiv, iar kilometrajul poate avea un coeficient negativ. \\
        KNN & Estimează prețul ca medie ponderată a celor mai apropiate anunțuri din setul de antrenare. Vecinii mai apropiați primesc pondere mai mare. & Un Golf din 2018 este comparat cu alte Golf-uri apropiate ca an, kilometraj și putere. \\
        Random Forest & Construiește mai mulți arbori de decizie pe eșantioane diferite și mediază predicțiile lor. & Un arbore poate învăța reguli pentru mașini diesel compacte, altul pentru mașini premium. \\
        Extra Trees & Seamănă cu Random Forest, dar folosește praguri de împărțire mai randomizate, crescând diversitatea arborilor. & Ajută când datele conțin multe interacțiuni și zgomot, dar poate fi mai puțin stabil pentru zone rare. \\
        Gradient Boosting & Construiește arbori secvențial. Fiecare arbore nou încearcă să corecteze erorile rămase de la ansamblul anterior. & Dacă modelul subestimează sistematic mașinile noi, arborii următori pot corecta parțial această zonă. \\
        Voting Ensemble & Combină predicțiile mai multor modele pentru a reduce dependența de o singură ipoteză. & Predicția finală folosește o medie ponderată pe baza performanței și, opțional, a acordului dintre modele. \\
        \bottomrule
    \end{tabular}
\end{table}

La nivel de implementare, toate modelele urmează același flux general:

\begin{lstlisting}[language=Python,caption={Pseudocod pentru antrenarea modelelor},label={lst:model-training-pseudo}]
date = incarca_csv()
X, y = separa_caracteristici_si_pret(date)
y_log = log1p(y)

pentru fiecare model candidat:
    pipeline = preprocesare + model
    pipeline.fit(X_train, y_log_train)
    predictii = expm1(pipeline.predict(X_test))
    calculeaza_RMSE_MAE_R2(predictii, y_test)

salveaza_modelul_cu_performanta_cea_mai_buna()
salveaza_metrici_grafice_si_predictii_individuale()
\end{lstlisting}

Această structură comună a permis compararea corectă a modelelor. Diferența dintre ele nu stă în datele primite sau în transformările aplicate, ci în modul în care fiecare model învață relația dintre caracteristici și preț.

\section{Antrenarea pe log-preț}
\label{sec:models-log-training}

Primele experimente au arătat că modelele pot produce predicții prea apropiate între anunțuri diferite, mai ales când distribuția țintei este dezechilibrată. Pentru a corecta parțial problema, modelele finale au fost antrenate pe \(log1p(price)\). Această transformare reduce presiunea exemplelor foarte scumpe și face ca erorile relative să conteze mai echilibrat.

În aplicație, predicțiile sunt inversate automat în EUR înainte de afișare. Utilizatorul nu vede valorile logaritmice. El vede prețul estimat, prețul anunțului și diferența procentuală.

\section{Predicția finală prin medie ponderată}
\label{sec:models-weighted}

Aplicația nu este limitată la un singur model. Pentru fiecare anunț, sunt calculate predicții pentru modelele salvate în bundle. Predicția finală afișată utilizatorului este o medie ponderată:

\[
    \hat{p}_{final} = \frac{\sum_i w_i \hat{p}_i}{\sum_i w_i}
\]

Prima strategie disponibilă, \texttt{inverse\_mae}, folosește ponderi invers proporționale cu MAE-ul modelului:

\[
    w_i = \frac{1}{MAE_i + \varepsilon}
\]

A doua strategie, \texttt{inverse\_mae\_with\_agreement}, adaugă o penalizare pentru modelele care se abat mult de la mediana predicțiilor:

\[
    w_i = \frac{1}{MAE_i + \varepsilon} \cdot exp\left(-\frac{|\hat{p}_i - median(\hat{p})|}{s}\right)
\]

Această a doua strategie a fost introdusă deoarece un model poate avea performanță bună global, dar poate produce o predicție foarte diferită pentru un anunț atipic. Penalizarea de acord nu elimină modelul, dar îi reduce influența în predicția finală.

\section{Verdictul}
\label{sec:models-verdict}

Verdictul compară prețul din anunț cu prețul estimat. Dacă diferența procentuală absolută este sub pragul ales de utilizator, anunțul este considerat corect. Dacă prețul publicat este mult sub estimare, verdictul devine preț prea mic, cu interpretare de ofertă potențial suspectă. Dacă prețul publicat este mult peste estimare, verdictul devine preț prea mare.

Pragul implicit este 15\%, dar interfața permite modificarea lui. Această opțiune este necesară deoarece toleranța diferă în funcție de buget și categorie. Pentru o mașină ieftină, 10\% poate fi o diferență importantă. Pentru o mașină premium, aceeași toleranță relativă poate corespunde unei sume mari, dar poate fi încă acceptabilă în negociere.


\section{Implementarea backend-ului}
\label{sec:implementation-backend}

Backend-ul pornește o aplicație FastAPI și montează folderul interfeței ca resurse statice. Modelele Pydantic validează automat tipurile și limitele parametrilor. De exemplu, toleranța trebuie să fie pozitivă și cel mult 100, iar metoda de agregare trebuie să fie una dintre opțiunile implementate.

\labelindexref{Fragmentul}{lst:backend-predict} arată partea centrală a endpoint-ului de predicție. Ruta nu conține logica matematică, ci rezolvă căile configurate, verifică artefactele și apelează serviciul de inferență. Această separare face posibilă folosirea aceleiași funcții și din CLI.

\begin{lstlisting}[language=Python,caption={Orchestrarea predicției în endpoint-ul FastAPI},label={lst:backend-predict}]
@app.post("/predict")
def predict(request: PredictRequest,
            user: AuthUser = Depends(get_current_user)):
    model_bundle = resolve_model_bundle_path()
    similarity_csv = resolve_similarity_csv_path()

    prediction = predict_from_link(
        site=request.site,
        url=request.url,
        model_bundle_path=model_bundle,
        threshold_percent=request.threshold_percent,
        similarity_csv_path=similarity_csv,
        ensemble_method=request.ensemble_method,
    )

    payload = prediction.to_dict()
    record_prediction_history(user_id=user.id,
                              prediction=payload)
    return payload
\end{lstlisting}

Funcția \texttt{predict\_from\_link} validează domeniul și selectează fluxul \autovit\ sau \mobilede. Rezultatul este reprezentat printr-o structură care conține atât verdictul, cât și datele intermediare necesare interfeței. Excepțiile de validare devin răspunsuri 400, iar erorile sursei externe sunt transformate în mesaje controlate.

\section{Implementarea frontend-ului}
\label{sec:implementation-frontend}

Frontend-ul folosește un formular HTML și o singură cerere asincronă către \texttt{/predict}. Înaintea trimiterii sunt validate domeniul linkului și toleranța. Pe durata cererii, butonul este dezactivat și este afișat overlay-ul de încărcare. Blocul \texttt{finally} este important deoarece închide animația atât la succes, cât și la eroare.

\begin{lstlisting}[language=Java,caption={Cererea frontend către API și închiderea stării de încărcare},label={lst:frontend-fetch}]
setLoading(true);
try {
  const response = await fetch("/predict", {
    method: "POST",
    credentials: "same-origin",
    headers: {"Content-Type": "application/json"},
    body: JSON.stringify({
      site,
      url,
      threshold_percent: thresholdPercent,
      ensemble_method: ensembleMethod,
      similar_limit: 4,
    }),
  });

  const data = await response.json();
  if (!response.ok) {
    throw new Error(data.detail);
  }
  renderResult(data);
} catch (error) {
  setMessage(error.message, "error");
} finally {
  setLoading(false);
}
\end{lstlisting}

Rezultatul este construit din carduri: informațiile vehiculului, verdictul, predicțiile pe model, graficele și anunțurile similare. Estimarea finală ponderată și verdictul sunt afișate imediat, iar estimările individuale și graficele sunt plasate într-un element HTML \texttt{details}, închis implicit. Această afișare progresivă păstrează concluzia ușor de citit și permite extinderea detaliilor numai când sunt necesare.

Frontend-ul tratează cele trei adrese \texttt{/}, \texttt{/istoric} și \texttt{/explicatii} ca pagini distincte, deși FastAPI servește același document HTML. JavaScript-ul activează secțiunea corespunzătoare adresei și marchează legătura curentă în bara de navigare. Pagina de explicații definește pe înțelesul utilizatorului metricile și regulile verdictului. Aceleași texte sunt disponibile printr-un element \texttt{dialog} deschis de butoanele de ajutor contextual.

Nu este folosit HTML primit de la platformele externe; valorile sunt introduse în șabloane controlate de aplicație și sunt escapate înainte de afișare. Preferința vizuală este memorată în \texttt{localStorage}, iar layout-ul folosește grid și flexbox pentru adaptarea cardurilor la desktop și mobil.

\labelindexref{Figura}{fig:gui-prediction} prezintă interfața după analiza unui anunț. În aceeași zonă sunt vizibile linkul, strategia de agregare, toleranța, confirmarea operației și cardul rezultatului. Fotografia ajută la verificarea identității anunțului, însă nu este folosită ca intrare pentru modelele curente.

\fig[width=0.98\textwidth]{src/img/carvaluator/gui_prediction.png}{fig:gui-prediction}{Interfața web după finalizarea unei predicții}

\section{Implementarea autentificării și securității}
\label{sec:implementation-security}

Parola este procesată cu PBKDF2-HMAC-SHA256. Pentru fiecare cont este generat un salt nou, iar baza păstrează algoritmul, numărul de iterații, saltul și hash-ul într-un șir serializat. Verificarea reconstruiește hash-ul și folosește \texttt{hmac.compare\_digest}, evitând comparațiile caracter cu caracter.

\begin{lstlisting}[language=Python,caption={Principiul de derivare sigură a parolei},label={lst:password-hash}]
salt = secrets.token_bytes(16)
derived = hashlib.pbkdf2_hmac(
    "sha256",
    password.encode("utf-8"),
    salt,
    PASSWORD_ITERATIONS,
)
password_hash = (
    f"pbkdf2_sha256${PASSWORD_ITERATIONS}$"
    f"{salt.hex()}${derived.hex()}"
)
\end{lstlisting}

Tokenul de sesiune este generat aleator și stocat în tabela \texttt{sessions}. Cookie-ul este HTTP-only, astfel încât JavaScript-ul paginii nu poate citi direct tokenul. Rutele de istoric primesc utilizatorul curent printr-o dependență FastAPI, iar comenzile SQL includ întotdeauna \texttt{user\_id}; un utilizator nu poate accesa sau șterge înregistrările altui cont. Endpoint-ul \texttt{GET /history/\{id\}} citește JSON-ul salvat și returnează raportul complet numai dacă perechea dintre identificator și utilizator este validă.

Validarea URL-urilor acceptă numai domeniile prevăzute. Aplicația nu execută cod din paginile analizate și nu expune în interfață conținut HTML brut. Pentru deployment, cookie-ul este marcat secure, iar toate cererile publice sunt transportate prin HTTPS-ul oferit de platformă.

\subsection{Limitarea ratei cererilor}

Limitatorul folosește câte o coadă cu momentele cererilor pentru fiecare pereche formată din scop și cheie. La o cerere nouă sunt eliminate momentele mai vechi decât fereastra configurată. Dacă numărul rămas a atins limita, cererea este respinsă înainte de hashing-ul parolei, scraping sau inferență. Accesul la structuri este protejat cu un lock deoarece endpoint-urile sincrone FastAPI pot rula concurent în thread-uri diferite.

\begin{lstlisting}[language=Python,caption={Verificarea unei limite cu fereastră glisantă},label={lst:rate-limit}]
decision = limiter.check(
    scope="predict",
    key=str(user.id),
    policy=settings.predict_policy,
)
if not decision.allowed:
    raise HTTPException(
        status_code=429,
        detail="Prea multe cereri de analiza.",
        headers={"Retry-After": str(decision.retry_after_seconds)},
    )
\end{lstlisting}

Configurația implicită permite 60 de cereri pe minut pentru un IP, 10 încercări de login în 15 minute, 3 înregistrări pe oră și 5 predicții pe minut pentru fiecare cont. Antetele \texttt{X-RateLimit-Limit}, \texttt{X-RateLimit-Remaining} și \texttt{Retry-After} fac comportamentul observabil. Ruta \texttt{/health} expune valorile active, fără a expune contoarele sau identitatea utilizatorilor.

\section{Implementarea similarității}
\label{sec:implementation-similarity}

Motorul de similaritate pornește de la aceeași reprezentare normalizată ca modelele. Sunt favorizate aceeași marcă și același model, apoi se calculează distanțe pentru an, kilometraj, putere și capacitate cilindrică. Combustibilul și transmisia adaugă bonusuri atunci când coincid.

Prețurile anunțurilor similare nu modifică predicția finală. Această decizie evită dublarea informației, deoarece modelele au fost deja antrenate pe aceleași tipuri de exemple. Lista are rol explicativ și de navigare, iar scorul său este separat de ponderile ensemble-ului.

Datasetul de similaritate este tratat ca un snapshot operațional distinct de setul de antrenare. Scriptul dedicat de actualizare colectează periodic anunțuri curente, le normalizează și actualizează fișierul stabil \texttt{similarity\_current.csv}. Separarea permite reducerea linkurilor expirate fără reantrenarea modelelor și fără modificarea estimărilor de preț.

\section{Artefacte și deployment pe Render}
\label{sec:implementation-deployment}

Antrenarea locală produce bundle-ul complet, metricile CSV, raportul JSON și graficele PNG. Scriptul PowerShell pentru pregătirea artefactelor Render creează o variantă redusă pentru cloud și copiază datasetul de similaritate. Fișierul \texttt{render.yaml} declară serviciul Docker, iar \texttt{Dockerfile} instalează Python 3.11, dependențele fixate și Chromium.

Fixarea versiunilor este necesară deoarece obiectele joblib folosesc clase interne din scikit-learn. Un model serializat cu versiunea 1.8.0 poate eșua la încărcare într-un runtime diferit. Containerul elimină această diferență și rezolvă și dependența Playwright de un browser instalat.

Render rulează aceeași aplicație, dar cu un bundle redus pentru limita de memorie a planului gratuit. Datasetul de similaritate este copiat separat în subdirectorul \texttt{datasets} al folderului \texttt{deploy\_artifacts}, sub numele \texttt{similarity\_current.csv}. Un refresh efectuat în folderul local \texttt{data} nu afectează automat site-ul public; artefactul trebuie regenerat, publicat în Git și urmat de redeploy. Accesul live la \mobilede\ rămâne dependent de politica platformei față de IP-urile de cloud; de aceea, lookup-ul în snapshot-ul livrat precedă scraping-ul live.
